% DK-08: first-page LLM-assistance disclosure.
%
% TMLR's FAQ requires authors who used LLM tools to say so "as a footnote on the
% first page", and states that "the ideas, claims and results found in
% submissions are human-sourced, not produced by such tools". The footnote below
% is therefore written to do two things: name the assistance accurately, and
% draw the line between assistance and authorship where TMLR draws it.
%
% DRAFT. It describes only the assistance Dingyi Kang used and can attest to.
% Chengshuai Yang and Ting Xue must each confirm their own use, or confirm they
% used none, before this is accurate as a statement for all three authors.
%
% Place immediately after \maketitle, which is where this was built and checked:
%     \maketitle
%     % DK-08: first-page LLM-assistance disclosure.
%
% TMLR's FAQ requires authors who used LLM tools to say so "as a footnote on the
% first page", and states that "the ideas, claims and results found in
% submissions are human-sourced, not produced by such tools". The footnote below
% is therefore written to do two things: name the assistance accurately, and
% draw the line between assistance and authorship where TMLR draws it.
%
% DRAFT. It describes only the assistance Dingyi Kang used and can attest to.
% Chengshuai Yang and Ting Xue must each confirm their own use, or confirm they
% used none, before this is accurate as a statement for all three authors.
%
% Place immediately after \maketitle, which is where this was built and checked:
%     \maketitle
%     % DK-08: first-page LLM-assistance disclosure.
%
% TMLR's FAQ requires authors who used LLM tools to say so "as a footnote on the
% first page", and states that "the ideas, claims and results found in
% submissions are human-sourced, not produced by such tools". The footnote below
% is therefore written to do two things: name the assistance accurately, and
% draw the line between assistance and authorship where TMLR draws it.
%
% DRAFT. It describes only the assistance Dingyi Kang used and can attest to.
% Chengshuai Yang and Ting Xue must each confirm their own use, or confirm they
% used none, before this is accurate as a statement for all three authors.
%
% Place immediately after \maketitle, which is where this was built and checked:
%     \maketitle
%     % DK-08: first-page LLM-assistance disclosure.
%
% TMLR's FAQ requires authors who used LLM tools to say so "as a footnote on the
% first page", and states that "the ideas, claims and results found in
% submissions are human-sourced, not produced by such tools". The footnote below
% is therefore written to do two things: name the assistance accurately, and
% draw the line between assistance and authorship where TMLR draws it.
%
% DRAFT. It describes only the assistance Dingyi Kang used and can attest to.
% Chengshuai Yang and Ting Xue must each confirm their own use, or confirm they
% used none, before this is accurate as a statement for all three authors.
%
% Place immediately after \maketitle, which is where this was built and checked:
%     \maketitle
%     \input{llm_assistance_disclosure}
% It then renders as an unnumbered footnote at the foot of page 1.

% An unnumbered footnote: \footnotetext[0] would print a literal "0" marker,
% and \thanks attaches to an author, which is wrong for a paper-level statement.
\begingroup
\renewcommand{\thefootnote}{}%
\footnotetext{%
\textbf{Use of LLM-based tools.}
The authors used an LLM-based coding assistant for engineering and verification
support during preparation of this submission: converting the manuscript to the
TMLR template, writing the analysis and regression-test scripts, independently
recomputing the reported values from the released evidence, auditing references
and internal consistency, drafting proposed wording corrections for the authors'
review, and packaging the reproducibility supplement.
The theory, the definitions, the experimental design, the measurements and the
results reported here are the authors' own and were not produced by these tools.
Every reported value was checked against the released evidence, and the authors
have verified and take full responsibility for the content of this paper.%
}%
\addtocounter{footnote}{-1}%
\endgroup

% It then renders as an unnumbered footnote at the foot of page 1.

% An unnumbered footnote: \footnotetext[0] would print a literal "0" marker,
% and \thanks attaches to an author, which is wrong for a paper-level statement.
\begingroup
\renewcommand{\thefootnote}{}%
\footnotetext{%
\textbf{Use of LLM-based tools.}
The authors used an LLM-based coding assistant for engineering and verification
support during preparation of this submission: converting the manuscript to the
TMLR template, writing the analysis and regression-test scripts, independently
recomputing the reported values from the released evidence, auditing references
and internal consistency, drafting proposed wording corrections for the authors'
review, and packaging the reproducibility supplement.
The theory, the definitions, the experimental design, the measurements and the
results reported here are the authors' own and were not produced by these tools.
Every reported value was checked against the released evidence, and the authors
have verified and take full responsibility for the content of this paper.%
}%
\addtocounter{footnote}{-1}%
\endgroup

% It then renders as an unnumbered footnote at the foot of page 1.

% An unnumbered footnote: \footnotetext[0] would print a literal "0" marker,
% and \thanks attaches to an author, which is wrong for a paper-level statement.
\begingroup
\renewcommand{\thefootnote}{}%
\footnotetext{%
\textbf{Use of LLM-based tools.}
The authors used an LLM-based coding assistant for engineering and verification
support during preparation of this submission: converting the manuscript to the
TMLR template, writing the analysis and regression-test scripts, independently
recomputing the reported values from the released evidence, auditing references
and internal consistency, drafting proposed wording corrections for the authors'
review, and packaging the reproducibility supplement.
The theory, the definitions, the experimental design, the measurements and the
results reported here are the authors' own and were not produced by these tools.
Every reported value was checked against the released evidence, and the authors
have verified and take full responsibility for the content of this paper.%
}%
\addtocounter{footnote}{-1}%
\endgroup

% It then renders as an unnumbered footnote at the foot of page 1.

% An unnumbered footnote: \footnotetext[0] would print a literal "0" marker,
% and \thanks attaches to an author, which is wrong for a paper-level statement.
\begingroup
\renewcommand{\thefootnote}{}%
\footnotetext{%
\textbf{Use of LLM-based tools.}
The authors used an LLM-based coding assistant for engineering and verification
support during preparation of this submission: converting the manuscript to the
TMLR template, writing the analysis and regression-test scripts, independently
recomputing the reported values from the released evidence, auditing references
and internal consistency, drafting proposed wording corrections for the authors'
review, and packaging the reproducibility supplement.
The theory, the definitions, the experimental design, the measurements and the
results reported here are the authors' own and were not produced by these tools.
Every reported value was checked against the released evidence, and the authors
have verified and take full responsibility for the content of this paper.%
}%
\addtocounter{footnote}{-1}%
\endgroup
